% =========================
% chuong/00-tom-tat.tex
% TÓM TẮT (Abstract)
% =========================

\section*{TÓM TẮT}
\addcontentsline{toc}{section}{TÓM TẮT}

Trong thương mại điện tử, tỉ lệ giao hàng đúng hạn thường được xem là thước đo chính của chất lượng
dịch vụ. Tuy nhiên, phân tích bộ dữ liệu công khai của Olist cho thấy 66,3\% đơn bị chấm 1 hoặc 2 sao
lại là đơn đã giao đúng hạn. Nhóm gọi hiện tượng này là sự không hài lòng ẩn và xây dựng một hệ thống
web giúp nhân viên chăm sóc khách hàng phát hiện và xử lý sớm những đơn có nguy cơ, trước khi khách
để lại đánh giá. Hệ thống nhận kết quả dự đoán rủi ro từ mô hình học máy, biến mỗi đơn có nguy cơ
thành một case trong hàng đợi ưu tiên, hỗ trợ nhận và giao việc theo giới hạn khối lượng, ghi nhận
từng lần liên lạc kèm minh chứng, đóng case với kết quả bắt buộc, và cho quản lý theo dõi tổng quan,
ra quyết định, cấu hình tham số, truy vết mọi thao tác. Nhóm phân tích nghiệp vụ bằng BPMN, FDD, DFD,
Use Case, thiết kế cơ sở dữ liệu PostgreSQL, xây dựng máy chủ API bằng FastAPI và giao diện chín màn
hình bằng HTML, CSS, JavaScript thuần, đóng gói bằng Docker Compose. Hệ thống được kiểm thử hai
đợt với dữ liệu 100 đơn Olist thật; các lỗi của đợt đầu đã được sửa và kiểm lại. Kết quả cho thấy có thể chuyển
một điểm số dự đoán thành quy trình chăm sóc khách hàng có trách nhiệm rõ ràng và truy vết được.

\noindent\textbf{Từ khóa:} thương mại điện tử, sự không hài lòng ẩn, chăm sóc khách hàng chủ động,
hệ thống hỗ trợ ra quyết định, ứng dụng web.
